% Paste from here into your main Overleaf document.
% Required packages (add to your preamble if not already there):
%   \usepackage{booktabs}
%   \usepackage{enumitem}

%%-----------------------------------------------------------------------
\chapter{Andre noter:}
\section{Research Question}

\textbf{How well can the Phillips Curve forecast inflation in Denmark compared
to the Eurozone, and has the relationship between unemployment and inflation
changed over time?}

The Phillips Curve posits a negative relationship between unemployment and
inflation. Denmark's unique labour market institutions (the ``flexicurity''
model) may produce a structurally different relationship from the broader
Eurozone. We exploit this cross-country contrast to apply the full toolbox
of the course, compare forecasting performance across models, and formally
test whether the Phillips Curve is stable over time.

%%-----------------------------------------------------------------------
\section{Data}



\textcolor{red}{
\subsection{Series}

\begin{table}[h!]
\centering
\begin{tabular}{llll}
\toprule
\textbf{Series} & \textbf{Frequency} & \textbf{Period} & \textbf{Source} \\
\midrule
Danish unemployment rate & Quarterly & 1990 onward & Statistics Denmark (dst.dk) \\
Danish HICP / CPI & Quarterly & 1990 onward & Statistics Denmark / Eurostat \\
Eurozone unemployment rate & Quarterly & 1995 onward & Eurostat \\
Eurozone HICP & Quarterly & 1995 onward & Eurostat / ECB \\
\bottomrule
\end{tabular}
\caption{Data overview}
\end{table}
}

\subsection{Transformations}

\begin{itemize}[nosep]
    \item Inflation is computed as the year-on-year percentage change in the
          price index.
    \item Unemployment is used as a rate in levels; stationarity will be
          assessed empirically.
    \item Log transformations and first differences are applied where unit root
          tests indicate non-stationarity.
    \item No population or calendar adjustments are expected to be necessary
          for these series, but this will be verified during exploratory
          analysis.
\end{itemize}

%%-----------------------------------------------------------------------
\section{Methods}

The analysis proceeds through eight stages, each building on the previous.
This structure ensures coverage of all major topics in the curriculum.

\subsection{Stage 1: Exploratory Analysis}

\begin{itemize}[nosep]
    \item Time plots for all four series, with visual inspection for trend,
          seasonality, and unusual observations.
    \item Seasonal subseries plots and STL decomposition to assess whether
          quarterly seasonality is present in unemployment.
    \item ACF and PACF plots to characterise autocorrelation structure.
    \item Augmented Dickey-Fuller (ADF) and KPSS unit root tests to determine
          the order of integration for each series.
\end{itemize}

\subsection{Stage 2: Structural Break Tests}

This is a central contribution of the paper. We formally test whether the
Phillips Curve relationship is stable or has shifted across macroeconomic
regimes.

\begin{itemize}[nosep]
    \item \textbf{Chow test} at candidate break dates: the 2008 financial
          crisis (2008 Q3), the COVID shock (2020 Q1), and the post-pandemic
          inflation surge (2021 Q3).
    \item \textbf{Bai-Perron multiple breakpoint test} to allow the data to
          determine the number and timing of breaks endogenously.
    \item \textbf{CUSUM and CUSUM-of-squares tests} to detect parameter
          instability over the full sample.
    \item Results will inform whether we estimate on the full sample or split
          the estimation window.
\end{itemize}

\subsection{Stage 3: General-to-Specific Model Selection}

Starting from a General Unrestricted Model (GUM), we test down to a
parsimonious specification following the Hendry/LSE methodology.

\begin{itemize}[nosep]
    \item Specify a general ADL model:
          \[
          \pi_t = \alpha + \sum_{i=1}^{p} \beta_i \pi_{t-i}
                  + \sum_{j=0}^{q} \gamma_j u_{t-j} + \varepsilon_t
          \]
          where $\pi_t$ is inflation and $u_t$ is the unemployment rate.
          Initial lag length is set generously (e.g.\ $p = q = 4$).
    \item Apply sequential $t$-tests and $F$-tests to eliminate insignificant
          lags, checking that residuals remain well-behaved at each step.
    \item Run this procedure separately for Denmark and the Eurozone and
          compare the resulting model specifications. Differences in retained
          lags reveal structural differences between the two labour markets.
    \item Information criteria (AIC, BIC) are used as a cross-check.
\end{itemize}

\subsection{Stage 4: Benchmark Models}

Simple models serve as a baseline for forecast comparison.

\begin{itemize}[nosep]
    \item Na\"{i}ve forecast (last observed value).
    \item Seasonal na\"{i}ve forecast.
    \item Simple exponential smoothing and Holt's method.
\end{itemize}

\subsection{Stage 5: Univariate ARIMA}

\begin{itemize}[nosep]
    \item Fit ARIMA$(p, d, q)$ models to the inflation series for Denmark and
          the Eurozone separately, using the Box-Jenkins identification
          procedure (ACF, PACF, unit root tests).
    \item Select the final specification via AIC / BIC.
    \item These models serve as a pure time-series benchmark: they exploit
          autocorrelation in inflation without reference to unemployment.
\end{itemize}

\subsection{Stage 6: Dynamic Regression (ARDL)}

\begin{itemize}[nosep]
    \item Estimate the parsimonious ADL model identified in Stage 3 as a
          dynamic regression model with unemployment as an external regressor.
    \item If structural breaks were found in Stage 2, include dummy variables
          for the relevant break dates or restrict estimation to the most
          recent stable regime.
    \item Inspect residuals for autocorrelation (Ljung-Box test), and
          heteroskedasticity (ARCH test). Address any remaining structure.
\end{itemize}

\subsection{Stage 7: Vector Autoregression (VAR)}

\begin{itemize}[nosep]
    \item Specify a bivariate VAR in $(\pi_t, u_t)$ for both regions.
    \item Select lag length via AIC / BIC / HQIC.
    \item Assess Granger causality: does unemployment Granger-cause inflation
          and vice versa?
    \item Generate impulse response functions (IRFs) to illustrate how an
          unemployment shock propagates to inflation over time.
    \item Produce multi-step-ahead forecasts from the VAR.
\end{itemize}

\subsection{Stage 8: Forecast Evaluation}

\begin{itemize}[nosep]
    \item Split the sample into an estimation window and a hold-out evaluation
          period (e.g.\ the final 8 quarters).
    \item Generate one-step-ahead and multi-step-ahead forecasts from all
          models using a rolling or expanding window.
    \item Compare accuracy using RMSE and MAE.
    \item Apply the \textbf{Diebold-Mariano test} to formally assess whether
          the best-performing model is statistically superior to the ARIMA
          benchmark.
    \item Report prediction intervals and assess calibration (are 95\%
          intervals covering roughly 95\% of actuals?).
\end{itemize}

%%-----------------------------------------------------------------------
\section{Paper Structure}

\begin{enumerate}[nosep]
    \item \textbf{Introduction} Research question, motivation, overview of
          approach.
    \item \textbf{Background} Phillips Curve theory, Denmark vs.\ Eurozone
          labour market institutions, prior empirical findings.
    \item \textbf{Data} Sources, transformations, descriptive statistics.
    \item \textbf{Exploratory Analysis} Time plots, decomposition, ACF/PACF,
          unit root tests.
    \item \textbf{Structural Break Analysis} Methodology and results.
    \item \textbf{Model Selection} General-to-specific procedure and final
          specifications for Denmark and the Eurozone.
    \item \textbf{Model Estimation} ARIMA, dynamic regression, VAR results.
    \item \textbf{Forecasting and Evaluation} Out-of-sample results,
          Diebold-Mariano test, prediction intervals.
    \item \textbf{Discussion} Cross-country comparison, what the structural
          differences reveal, limitations.
    \item \textbf{Conclusion} Main findings and directions for future work.
\end{enumerate}

%%-----------------------------------------------------------------------
\section{Division of Work}

\begin{table}[h!]
\centering
\begin{tabular}{ll}
\toprule
\textbf{Section} & \textbf{Responsible} \\
\midrule
Data collection and cleaning & TBD \\
Exploratory analysis (Stage 1) & TBD \\
Structural break tests (Stage 2) & TBD \\
General-to-specific (Stage 3) & TBD \\
ARIMA and benchmarks (Stages 4 and 5) & TBD \\
Dynamic regression (Stage 6) & TBD \\
VAR (Stage 7) & TBD \\
Forecast evaluation (Stage 8) & TBD \\
Writing and editing & All \\
\bottomrule
\end{tabular}
\caption{Work distribution (to be completed by the group)}
\end{table}

%%-----------------------------------------------------------------------
\section{Software and Replication}

All analysis will be conducted in \textbf{R}. Key packages include:
\texttt{forecast}, \texttt{fable}, \texttt{tseries}, \texttt{strucchange}
(structural breaks), \texttt{vars} (VAR models), and \texttt{dynlm} (dynamic
regression). Code will be shared via a joint GitHub repository.

% End of paste